\section{Security Analysis and Discussion}
\label{sec:discussion}

\subsection{What the Evidence Shows}

The two studies answer different questions.  The independent attack
reproduction shows that a bounded SABER instruction perturbation can propagate
through an unprotected \vla to the registered physical-risk endpoint.  Its
observed ASR is 45.35\% among 86 clean-eligible units (39/86).  The four-arm
study then characterizes the final \system prototype and its stage ablations at
the consumer boundary.  It shows predispatch rejection for the covered C1 findings, the
expected response to the exercised C2 transaction faults, and 0.00\%
trace-level joint-limit incidence across 237 valid attacked \ltwoenabled traces.

The system-level result does not show a broad safety improvement.  On the common
75-unit cohort, the broad risk-transition rates are 50.67\%, 56.00\%, 48.00\%,
and 57.33\%, and no monitored arm provides evidence of a reduction relative to
\vlaonly.  Under attack, \dual task success is 53.33\%, compared with 60.83\%
for \vlaonly.  Moreover, each condition is 98.96\% valid against a registered
100\% requirement, so the four-arm results cannot support a confirmatory
overall-effect claim.  The
paper therefore treats the C1, C2, and joint-side observations as bounded
property evidence, not as proof that \system makes the robot generally safer.

\subsection{Why Local Properties Do Not Imply End-to-End Safety}

The evaluation illustrates why the three components must be interpreted
separately.  \lone and \ltwo protect different alignment stages; \ltwoa and
\ltwob supply distinct evidence within the second stage.  \taskassessment
enforces a declared checker vocabulary against an
independently anchored task state.  It can stop a covered hard finding before
dispatch, but an incomplete checker can reject useful actions or miss hazards
outside its predicates.  The observed task-completion loss is therefore not an
incidental reporting detail; it is the availability cost of the current checker
boundary.

\transactionbinding answers a different question.  It determines whether
execution evidence still refers to the proposal that received permission.  It prevents a covered
substitution or replay from being mistaken for aligned execution, but a faithfully
executed authorized action can still be physically undesirable.  It also cannot
detect a trusted sink that changes both the action and the evidence it reports.

\guardscreening targets one registered joint-side condition.  Its zero
joint-limit observation in valid attacked \ltwoenabled traces is consistent with that
local goal, while the broader endpoint still includes contact, force,
cost/collision, and other pathways not eliminated by the joint guard.  Changing
the guard/controller configuration can also alter task progress even though the
policy-emitted source action remains fixed.  Consequently, source-action
integrity, local constraint enforcement, task success, and broad physical risk
are related but non-equivalent properties.

\subsection{Dependence on the TCB and Registered Models}

The authorization claim depends on a task and observation path that is
independent of the attacked policy view.  The evaluated prototype realizes this
path with privileged simulator task state and geometry.  A deployable camera-only
or sensor-based trusted path would need its own integrity and accuracy evidence;
errors before the trusted split are outside the current guarantee.

The execution claim assumes complete mediation, protected monitor state, and
truthful receipts and actuator feedback.  Contract binding ranges only over the
effects the contract declares, so omitted hazards remain uncovered.  Likewise,
the checker-relative claim depends on the registered predicates.  Unknown
inputs trigger configured fail-closed handling; no general semantic judgment
follows.

The guard evidence uses the LIBERO/robosuite dynamics
~\cite{robosuite2020} and the same simulator model for shadow screening and
post-restore execution.  It qualifies snapshot restoration and the registered
joint-margin and force measurements in that environment.  It does not establish
model fidelity, robustness to real sensor noise, or physical-robot safety.  The
Lean model is narrower still: it checks selected abstract authorization,
ordered-action, evidence, and phase relations, not the prototype-to-model
refinement or continuous dynamics.

\subsection{Statistical and Empirical Limitations}

The four-arm study produced 1.04\% invalid rows in each condition (5/480)
against a registered zero-invalid requirement.  Task counts retain all 120 units with
invalid runs counted as failures, but the broad-risk comparison uses 75 units
that are eligible and valid in every arm.  That common cohort enables a matched
comparison with a shared denominator; it is not a full-population causal
estimate.  The zero joint-limit result similarly covers only the 237 valid
attacked \ltwoenabled traces and cannot be extended to invalid runs.

The empirical scope contains one $\pi_{0.5}$ checkpoint, one simulator
benchmark, and one frozen SABER instruction-attack family.  Defense-aware
adaptive attacks are not evaluated, although recent work shows their importance
~\cite{adaptivesecond2026}.  Other visual or post-inference channels belong to
the threat model but are represented here only by focused transaction cases, not
full end-to-end attack campaigns.  Finally, the 458.99\,ms maximum screening
latency violates the registered 200\,ms bound, so the prototype has no
hard-real-time guarantee.  Broader attack families, trusted perception, more
complete physical predicates, and real-robot evaluation are necessary before
claiming wider protection.
